% To follow abelian_components.tex.  Uses its notation and proved lemmas.
% No private macros beyond the main paper's existing conventions.

\section{A canonical nonabelian quotient of class two}
\label{sec:class-two-components}

We construct a canonical class-two quotient for every odd prime.
Its full odd commutator module is detected by the outer action;
the cup comparison determines which even abelian components survive.

Retain the notation $W=U/D$ and the class-two Lie module constructed
in Section~\ref{sec:abelian}.  Write it as
\[
 \mathcal L=W\oplus\mathcal E,\qquad
 0\longrightarrow\widehat{\bigwedge}^{\,2}W
 \longrightarrow\mathcal E\longrightarrow T\longrightarrow0.
\]
The first summand is its canonical $c$-negative part, the second is
central and $c$-positive, and $[w,w']=w\wedge w'$ in the displayed
central submodule.  The direct sum $W\oplus\mathcal E$ is an
identity of additive Lie modules; no section of
$\mathcal E\to T$ is chosen.

For a set $E$ of even characters define the central Lie ideal
and its quotient by
\[
 \mathcal I_E=\bigoplus_{b\notin E}\mathcal E^{[b]},\qquad
 \mathcal L_E=\mathcal L/\mathcal I_E,
 \qquad C_E=(\mathcal L_E,\star),\qquad
 x\star y=x+y+\tfrac12[x,y].
\]
Only even $b$ occur in the direct sum.  These are compact modules,
and $2$ is a unit in $\mathbb Z_p$, so $C_E$ is a pro-$p$ group.
Let $q_E:H\twoheadrightarrow C_E$ be the composite of the
class-two quotient, the quotient killing the canonical odd section
of $D$ and its brackets, and the quotient by $\mathcal I_E$.
In additive coordinates its bracket is
\[
 [w,w']_{C_E}=\operatorname{pr}_E(w\wedge w').
\]
Thus discarding $b$ removes both $T_b$ and the central brackets
in character $b$.  This formula specifies every map used in the
construction.

Put $Z=\mathcal E_{\Lambda\text{-tor}}$ and
$Q=(W\oplus\mathcal E/Z,\star)$. The proof below constructs
an injection of $Q$ into a quotient of the rational outer image.


\begin{theorem}[Canonical class-two detection]
\label{thm:canonical-torsion-quotient}
For every odd prime $p$, the group $Z$ is closed, finite free
 over $\mathbb Z_p$, and saturated in $\mathcal E$; it has
zero intersection with $N=\widehat{\bigwedge}^{\,2}W$.
The kernel of $q$ is normal in $\operatorname{Gal}(\Omega_p/\mathbb Q)$
and contains the pro-$p$ outer-action kernel in $H$. Thus
\[
 L^{(2)}=\Omega_p^{\ker q}\subseteq\Omega_p^*,\qquad
 \operatorname{Gal}(L^{(2)}/F_\infty)=Q.
\]
The group $Q$ is torsionfree and nonabelian of class two, with
\[
 [Q,Q]=N,\qquad Q^{\rm ab}=W\oplus T/\pi_T(Z),
\]
where $\pi_T:\mathcal E\to T$ is the displayed projection.
For every even $e$ such that $X^{[e]}$ is finite,
$Z^{[e]}=0$; consequently the full $T_e$ occurs in the even
abelianization of $Q$.
\end{theorem}

We first record the separation of the completed commutator module
used in the proof.  It is separation by generator length in the
full free Lie algebra of the rational outer image.

\begin{lemma}\label{lem:class-two-moments}
Let $\mathcal M$ be the odd weights retained after removing the
finite exceptional set in Section~\ref{sec:abelian}.
Write the linear coefficients of that rational outer
representation as
\[
 \kappa_m(w)=b_m F_a(u^m-1),\qquad
 a\equiv m\pmod{p-1},\quad b_m\ne0.
\]
The induced bracket map from
$\widehat{\bigwedge}^{\,2}_{\mathbb Z_p}W$
to the length-two part of the weight-completed free Lie algebra
on the $s_m$, $m\in\mathcal M$, is injective.
The assertion also holds after projecting source and target onto
any collection of $\Delta$-characters.
\end{lemma}

\begin{proof}
Choose the rational $R$-basis of $W[1/p]$ from
\eqref{eq:Wfree}.  Clearing denominators in this basis and its
inverse embeds $W$ in a finite sum of copies of $\Lambda$,
with cokernel killed by a fixed power of $p$.
The same construction embeds its completed tensor square
in a finite sum of $\mathbb Z_p[[T_1,T_2]][1/p]$.
For precision, a reverse map multiplied by $p^r$ gives a
composite equal to $p^r$ on $W$; tensoring twice gives
$p^{2r}$.  The completed tensor square is $p$-torsionfree,
as is seen from a product $\mathbb Z_p$-basis of $W$, so
the tensor map is injective.  Since $2$ is invertible,
the completed exterior square is its alternating summand.

Let $A_{a,b}(T_1,T_2)$ be its alternating coordinates.
The coefficient at $[s_m,s_n]$ is
\[
             b_mb_n A_{a,b}(u^m-1,u^n-1).
\]
Use antisymmetry to read a coefficient with $m>n$ as the
negative of that with the two indices reversed.  Thus these
coefficients are available for every distinct ordered pair,
not only for $m<n$.
If they all vanish, fix $n$ and vary $m$ through its residue
class in $\mathcal M$, excluding $m=n$.  This is an infinite
set of distinct points in $p\mathbb Z_p$.
The one-variable Weierstrass theorem makes
$A_{a,b}(T_1,u^n-1)$ identically zero.  Varying $n$ and
applying the same argument to each coefficient of $T_1$
gives $A_{a,b}=0$.  Every coordinate vanishes.
For a fixed character $e$, only pairs $a+b=e$ occur;
the same argument uses exactly those pairs and proves
the projected assertion.
\end{proof}

\begin{lemma}\label{lem:class-two-invariance}
The quotient $\mathcal L$ and each central character quotient $C_E$ have kernels stable under $\operatorname{Gal}(\Omega_p/\mathbb Q)$.
\end{lemma}
\begin{proof}
We first check the invariance and the canonical quotient map.
Before removing $D$, the first quotient can be written
intrinsically as
\[
 H/\overline{[\pi^{-1}(T),H]},\qquad \pi:H\longrightarrow M.
\]
Indeed $[H,H]\subseteq\pi^{-1}(T)$, so its third lower
central subgroup is killed, and the further killed central
brackets are precisely those involving $T$.
Its central subgroup above $T$ is the $c$-positive Lie
module $\mathcal E$ before removing $[D,U]$.
The ideal generated by the canonical negative lift of $D$
is that lift plus $[D,U]$.
It is preserved by $\Gamma$ and $\Delta$, and it is an ideal,
so it is preserved by conjugation from $H$ as well.
Here is the change-of-section calculation for full arithmetic
invariance. For $g\in\operatorname{Gal}(\Omega_p/\mathbb Q)$ write
$gcg^{-1}=\operatorname{Ad}(h)c$ and let $v$ be the odd abelian
part of $\log h$. If $s_c:U\to\mathcal L_0^-$ is the negative
section before removing $D$, write
$g s_c(u)=s_c(g_Uu)+\lambda_g(u)$ with central $\lambda_g(u)$.
The identity $gc=\operatorname{Ad}(h)cg$ gives
\[
 -s_c(g_Uu)-\lambda_g(u)
 =-s_c(g_Uu)+\lambda_g(u)-[v,g_Uu],
 \qquad \lambda_g(u)=\tfrac12[v,g_Uu].
\]
As $g_U(D)=D$, the difference belongs to $[D,U]$. Thus the
ideal is invariant under every $g$, independently of the section.
The action on the central positive part is unaffected by inner
conjugations.  Consequently
its character summands, and hence $\mathcal I_E$, are
independent of these lift choices.  The closed subgroup
$\ker q_E$ is normal in $\operatorname{Gal}(\Omega_p/\mathbb Q)$.

\end{proof}

\begin{proof}[Proof of Theorem~\ref{thm:canonical-torsion-quotient}]
We first establish the module assertions used to form the quotient.
In the bounded coordinates for $N$ constructed above, the action
of $\gamma$ is multiplication by $(1+T_1)(1+T_2)$.
For any nonzero $f\in\mathbb Z_p[[\gamma-1]]$, the series
\[
             f((1+T_1)(1+T_2)-1)
\]
is nonzero: setting $T_2=0$ recovers $f(T_1)$.
The coefficient rings $\mathbb Z_p[[T_1,T_2]][1/p]$ are
integral domains. Multiplication by every nonzero $f$ is
therefore injective on $N$, so $N$ is $\Lambda$-torsionfree.

The module $T$ is finite free over $\mathbb Z_p$.
Let $a(z)$ be the characteristic polynomial of $\gamma$ on
$T$; if $T=0$, put $a=1$. Cayley--Hamilton gives
$a(\gamma)\mathcal E\subseteq N$. If $x\in Z$, then
$a(\gamma)x$ is a $\Lambda$-torsion element of $N$, hence zero.
Conversely every element killed by $a(\gamma)$ is
$\Lambda$-torsion. Thus
\[
                     Z=\ker(a(\gamma):\mathcal E\to\mathcal E).
\]
This identifies $Z$ as a closed submodule. Its intersection
with $N$ is zero, so $\pi_T$ embeds $Z$ in the finite free
$\mathbb Z_p$-module $T$. In particular $Z$ is finite free
over $\mathbb Z_p$. The module $\mathcal E$ is
$\mathbb Z_p$-torsionfree, and the kernel of an endomorphism
of such a module is saturated: $px\in Z$ implies
$p\,a(\gamma)x=0$, hence $a(\gamma)x=0$ and $x\in Z$.
Consequently $\mathcal E/Z$ is also torsionfree.

The arithmetic action on the central module $\mathcal E$
factors through $\Delta\times\Gamma$. Its torsion submodule
$Z$ is invariant under that action. The preceding intrinsic
construction of $\mathcal L$, including the change-of-section
calculation, proves invariance under the entire arithmetic
group. Thus $q$ has Galois-stable kernel.

Take the rational free outer image constructed above. In
addition to the finite exceptional odd weights coming from
$D$, remove every odd letter occurring in a positive pair
$m,n\geq3$ with
\[
                              a(u^{m+n})=0.
\]
There are only finitely many such pairs: $a$ has finitely
many roots, $u$ has infinite order, and any fixed positive
sum has finitely many positive summands. The remaining
set of odd weights is cofinite. Density in each remaining
abelian weight quotient still makes its linear coefficient
$b_m$ nonzero.

Pass to generator length at most two in this free target.
The linear image of $T$ is zero by parity, so its brackets
have length at least three. The canonical odd section of $D$
has zero linear image by the choice of weights and zero
quadratic image by $c$-parity. Its brackets also vanish.
Thus the map factors through $\mathcal L$.
The image of $Z$ is central of length two. Its coefficient
at $[s_m,s_n]$ is killed by $a(u^{m+n})$, which is nonzero
by construction. Every such coefficient is therefore zero.
We obtain a factorization on the whole group
\[
 H\xrightarrow{q} Q\xrightarrow{j}
      \exp(\mathfrak f_{\leq2}),
\]
where $\mathfrak f$ is free on the remaining odd letters.

We prove that $j$ is injective. The degree-one moment
calculation detects $W$, so an element in its kernel lies
in $\mathcal E/Z$. Represent it by $x\in\mathcal E$.
Then $a(\gamma)x\in N$ has zero image in the length-two
target. The bracket moment lemma is still valid after
deleting any finite set of odd weights, and is injective
on all of $N$. Hence $a(\gamma)x=0$. The kernel identity
for $Z$ proved above gives $x\in Z$, as required.

The pro-$p$ outer kernel maps to zero in this rational
quotient. The factorization and injectivity therefore give
its containment in $\ker q$. Galois correspondence proves
the field assertions. Since $Z\cap N=0$, the commutator
module remains exactly $N$; the abelianization is
$W\oplus T/\pi_T(Z)$. The additive Lie module
$W\oplus\mathcal E/Z$ is torsionfree, and powers in a
class-two BCH group satisfy $x^{p^r}=p^rx$. Thus $Q$ has
no nontrivial finite-order element. Its commutator module
is nonzero, because $W[1/p]$ has infinite dimension.

Finally fix an even $e$ with $X^{[e]}$ finite.
Suppose $Z^{[e]}[1/p]$ is nonzero. Over a finite splitting
field $L/\mathbb Q_p$ it contains an eigenvector $x$ of eigenvalue $\alpha$.
Its nonzero image in $T_e$ belongs to the unique
$\alpha$-eigenline of that cyclic module. In the notation
of the companion identity, there is a nonzero scalar $d$
with
\[
 \pi_T(x)=d\,r(\gamma)v_e,\qquad
 r(z)=q_e(z)/(z-\alpha).
\]
Write $x=d\,r(\gamma)y_e-A$, with $A\in N\otimes_{\mathbb Z_p}L$.
Applying $\gamma-\alpha$ gives
\[
                         d\,C_e=(\gamma-\alpha)A.
\]
Every coefficient of $C_e$ would therefore vanish on
$(1+T_1)(1+T_2)=\alpha$, contradicting the proved cup-pole
lemma. Hence $Z^{[e]}=0$. For a set $E$ of such characters,
$Z\subseteq\bigoplus_{b\notin E}\mathcal E^{[b]}$.
Thus the original map $q_E$ factors through $q$, and
$\ker q\subseteq\ker q_E$. The corresponding field
inclusion is $L_E^{(2)}\subseteq L^{(2)}$.
\end{proof}



\begin{corollary}[Canonical class-two component detection]
\label{thm:class-two-components}
Let $p$ be odd, and let $E$ be a set of even exponents for which
$X^{[e]}$ is finite for every $e\in E$.
Then the pro-$p$ outer-action kernel in $H$ is contained
in $\ker q_E$.  The field
\[
 L_E^{(2)}=\Omega_p^{\ker q_E}
\]
is Galois over $\mathbb Q$ and is contained in $\Omega_p^*$.
Its Galois group over $F_\infty$ is $C_E$, and
\begin{align*}
 C_E^{\rm ab}&=(U/D)\oplus\bigoplus_{e\in E}T_e=B_E,\\
 [C_E,C_E]&=\bigoplus_{e\in E}
       \left(\widehat{\bigwedge}^{\,2}_{\mathbb Z_p}(U/D)\right)^{[e]}.
\end{align*}
The largest subextension of $L_E^{(2)}/F_\infty$ which is
abelian over $F_\infty$ is the field $L_E$ in
Corollary~\ref{thm:abelian}.  For nonempty $E$, the group $C_E$
is nonabelian.
\end{corollary}

\begin{proof}[Proof of Corollary~\ref{thm:class-two-components}]
For every selected $e$, Theorem~\ref{thm:canonical-torsion-quotient}
gives $Z^{[e]}=0$. Hence $q_E$ factors through $q$, and
$L_E^{(2)}\subseteq L^{(2)}\subseteq\Omega_p^*$.
Lemma~\ref{lem:class-two-invariance} gives normality over
$\mathbb Q$. The displayed bracket formula gives the
commutator and abelianization, including the maximal abelian
subfield assertion.

Finally each odd character of $W[1/p]$ is a nonzero rank-one
$R$-module and has infinite dimension over $\mathbb Q_p$.
For any even $e$, choose odd $a,b$ with $a+b=e$.
For $a\ne b$, nonzero vectors in these two summands have
nonzero exterior product.  For $a=b$, choose two
$\mathbb Q_p$-independent vectors in that summand.
Clearing denominators places them in $W$ and preserves their
nonzero wedge.  Thus every selected central bracket
component is nonzero, proving nonabelianness for nonempty $E$.
\end{proof}




\begin{corollary}[Abelian component detection]\label{thm:abelian}
For the same set $E$, put $B_E=W\oplus\bigoplus_{e\in E}T_e$
and let $L_E$ be the fixed field of the kernel of the canonical
map $H\to H^{\rm ab}\to B_E$. Then $L_E/\mathbb Q$ is Galois,
$\operatorname{Gal}(L_E/F_\infty)=B_E$, and
$L_E\subseteq L_E^{(2)}\subseteq\Omega_p^*$.
\end{corollary}
\begin{proof}
The abelianization of $q_E$ is exactly the displayed map to $B_E$.
Thus $\ker q_E$ is contained in its kernel. Galois correspondence
and Corollary~\ref{thm:class-two-components} give both inclusions.
The kernel is stable under the full arithmetic group because the
character projections and $D$ are stable.
\end{proof}
